\documentclass[11pt]{article}
\usepackage[margin=1in]{geometry}
\usepackage{graphicx}
\usepackage{booktabs}
\usepackage{longtable}
\usepackage{array}
\usepackage{tabularx}
\usepackage{multirow}
\usepackage{float}
\usepackage{caption}
\usepackage{subcaption}
\usepackage{amsmath}
\usepackage{amssymb}
\usepackage{siunitx}
\usepackage{url}
\usepackage[hidelinks]{hyperref}
\usepackage{xcolor}
\usepackage{enumitem}
\usepackage{pdflscape}

\title{Regional Background-Rate Assessment of the 2025--2026 Aomori Seismic Activity}
\author{TRACE}
\date{\today}


% --- EarthLink enforced font settings (XeLaTeX) ---
\usepackage{fontspec}
\usepackage{xeCJK}
\usepackage{unicode-math}
\setmainfont{DejaVu Serif}
\setCJKmainfont{Noto Serif CJK SC}
\setmathfont{XITS Math}
\usepackage[htt]{hyphenat}
% --- end enforced font settings ---

\begin{document}
\maketitle

\begin{abstract}
This report evaluates whether the 2025--2026 Aomori activity was exceptional relative to the 2020--2026 background seismicity, using the long-term raw JMA/Hi-net catalog as the background reference and a relocated/filtered active-period catalog as the high-resolution sequence dataset. To preserve catalog consistency, the long-term catalog was filtered to a common spatial domain defined from the relocated active-period footprint before any background comparison. The common region spans 38.353003--42.532633$^{\circ}$N and 140.851921--144.64847$^{\circ}$E. Cross-catalog audit results show that larger events are highly consistent in the overlap period and support robust long-term comparison primarily for $M\geq 3$, $M\geq 4$, and $M\geq 5$.

Within that harmonized domain, the 2025-10 to 2026-05 active period is clearly exceptional relative to the 2020--2026 background. In the common region, observed versus expected counts are 1267 versus 348.38 for $M\geq 3$, 259 versus 62.62 for $M\geq 4$, 68 versus 12.03 for $M\geq 5$, and 13 versus 1.95 for $M\geq 6$, corresponding to percentile ranks of about 98.5--99.6. The strongest anomaly is concentrated in the shallow 0--30 km portion of the \texttt{m1\_m3\_local\_union}, which remains extreme after bootstrap and random-window controls. The anomaly is therefore best interpreted as a localized-to-subregional rate pulse superposed on a weaker broader regional elevation, rather than as a spatially uniform regional increase. The results are statistical only and do not demonstrate physical triggering or any specific mechanism.
\end{abstract}

\section{Scientific Objective}
The objective was to build a regional background-rate model for the Aomori earthquake catalog system and use it to test whether the 2025--2026 activity was exceptional relative to the 2020--2026 background. The required workflow separated two catalog roles:
\begin{itemize}[leftmargin=1.5em]
    \item the long-term raw catalog as the background-rate reference; and
    \item the relocated/filtered active-period catalog as the fine-scale dataset for spatial, temporal, and depth decomposition during the active interval.
\end{itemize}
A critical constraint was that the two catalogs were \emph{not} to be merged as if they were homogeneous products. Instead, a common analysis region had to be defined from the active relocated catalog footprint, and the long-term raw catalog had to be filtered to that same spatial mask before any long-term versus active-period comparison. This report follows that requirement explicitly.

The main scientific questions were:
\begin{enumerate}[leftmargin=1.5em]
    \item Is the 2025--2026 Aomori active period exceptional relative to the 2020--2026 background within the common analysis region?
    \item Which regions and depth ranges show the strongest rate anomalies?
    \item Are the anomalies localized or part of a broader regional rate pulse?
    \item Does the M1--M3 local region remain anomalous after background correction?
    \item Can M2 outer-band activity be explained by broader regional background changes?
    \item Which observations are statistically strong enough to justify later physical follow-up?
\end{enumerate}

\section{Data, Catalog Roles, and Implementation Choices}
\subsection{Catalogs and actual data used}
The long-term reference catalog was the raw product spanning 2020-01-01 to 2026-05-22. The requested active-period file was unavailable, and the workflow substituted the resolved relocated catalog documented in the task diagnostics. The actual active-period catalog used in the analysis was:
\url{<CASE_ROOT>/data/catalog/Snet_catalog_relocate_250930_260501.csv}.
This substitution is documented in the task outputs and should be retained in any reproducibility discussion.

\subsection{Common-region definition}
A common analysis region was defined from the active relocated catalog footprint with a 0.15$^{\circ}$ buffer, and the long-term raw catalog was filtered to that same mask before background-rate estimation. The resulting common region is:
\begin{center}
38.353003--42.532633$^{\circ}$N, \quad 140.851921--144.64847$^{\circ}$E.
\end{center}
This step is essential because the long-term raw catalog extends well beyond the relocated footprint and would otherwise bias rate comparisons.

Figure~\ref{fig:footprints} shows the footprint mismatch and the adopted common domain.

\begin{figure}[H]
    \centering
    \includegraphics[width=0.82\textwidth]{<CASE_ROOT>/run/02_4_bg_rate_correction/exp_run/outputs/01_aomori_background_rate_analysis/figures/01_catalog_footprints_common_region.png}
    \caption{Catalog footprints and the common analysis region used for all primary long-term versus active-period comparisons. The long-term raw catalog extends beyond the relocated active-period footprint, which justifies spatial harmonization before anomaly testing.}
    \label{fig:footprints}
\end{figure}

\subsection{Catalog audit and threshold reliability}
The audit established the evidence chain needed for defensible rate comparison:
\begin{itemize}[leftmargin=1.5em]
    \item Long-term raw catalog size: 227,755 events.
    \item Active relocated catalog size: 22,096 events.
    \item After common-region filtering: 159,087 long-term events and 22,090 active-period events.
    \item Retention fractions: 0.6985 for the long-term catalog and 0.9997 for the active catalog.
    \item Estimated magnitude completeness by maximum curvature: $M_c = 1.15$ for the long-term common-region catalog and $M_c = 1.45$ for the active common-region catalog.
\end{itemize}
These results support the following threshold policy:
\begin{itemize}[leftmargin=1.5em]
    \item \textbf{Preferred for cross-catalog background comparison:} $M\geq 3$, $M\geq 4$, $M\geq 5$.
    \item \textbf{Descriptive only:} $M\geq 2$ where count support is limited.
    \item \textbf{Restricted to active-period internal analysis:} $M\geq 1.2$, unless additional completeness justification is provided.
\end{itemize}

\subsection{Crosswalk consistency in the overlap period}
Within the overlap period and common region, event counts are similar at larger magnitudes: raw versus relocated counts are 1389 versus 1267 for $M\geq 3$, 280 versus 259 for $M\geq 4$, 69 versus 68 for $M\geq 5$, and 13 versus 13 for $M\geq 6$. Matched-event diagnostics show essentially identical magnitudes and very small differences in origin time, location, and depth. Figure~\ref{fig:matchdiff} documents this consistency and supports the use of higher-magnitude thresholds for the long-term background comparison.

\begin{figure}[H]
    \centering
    \includegraphics[width=0.82\textwidth]{<CASE_ROOT>/run/02_4_bg_rate_correction/exp_run/outputs/01_aomori_background_rate_analysis/figures/03_matched_event_differences.png}
    \caption{Matched large-event differences between the long-term raw and relocated catalogs in the overlap period. The small timing, location, and depth offsets support strong cross-catalog consistency for larger events.}
    \label{fig:matchdiff}
\end{figure}

\section{Background-Rate Modeling Strategy}
The long-term background model was estimated from the spatially filtered 2020--2026 raw catalog, using monthly windows and region-specific baselines. The active-period decomposition used the relocated catalog for internal fine-scale temporal, regional, and depth structure. This separation preserves the intended roles of each catalog:
\begin{enumerate}[leftmargin=1.5em]
    \item the raw catalog provides the background-rate reference; and
    \item the relocated catalog provides the active-period structural detail.
\end{enumerate}
The primary regions discussed in the evidence are the \texttt{common\_region}, \texttt{m1\_m3\_local\_union}, \texttt{m2\_near\_field}, \texttt{m2\_outer\_band}, and \texttt{control\_region}. Depth partitions include 0--30 km, 30--60 km, and $>60$ km. Heavy bootstrap/random-window controls were applied only to $M\geq 3$, $M\geq 4$, and $M\geq 5$, with 4000 resamples per combination.

\section{Results}
\subsection{Was the 2025--2026 active period exceptional relative to the 2020--2026 background?}
Yes. Within the common region, the full active period from 2025-10 to 2026-05 is clearly exceptional relative to the long-term background. The main quantitative results are:
\begin{center}
\begin{tabular}{lcccc}
\toprule
Threshold & Observed & Expected & Obs./Exp. & Percentile \\
\midrule
$M\geq 3$ & 1267 & 348.38 & 3.64 & 98.49 \\
$M\geq 4$ & 259 & 62.62 & 4.14 & 98.49 \\
$M\geq 5$ & 68 & 12.03 & 5.65 & 98.73 \\
$M\geq 6$ & 13 & 1.95 & 6.68 & 99.62 \\
\bottomrule
\end{tabular}
\end{center}
These statistics come from the primary background-comparison outputs of Task 01. They show that the active-period counts exceed expectation by factors of about 3.6 to 6.7 across robust thresholds, with percentile ranks near the top of the 2020--2026 distribution.

The bootstrap and random-window control analysis further strengthens this conclusion for the robust thresholds $M\geq 3$, $M\geq 4$, and $M\geq 5$. For the shallow common region (0--30 km), the control summary reports:
\begin{itemize}[leftmargin=1.5em]
    \item $M\geq 3$: primary obs./exp. $= 5.38$, random percentile $= 98.8$, outside-period percentile $= 100.0$;
    \item $M\geq 4$: primary obs./exp. $= 5.82$, random percentile $= 98.8$, outside-period percentile $= 100.0$;
    \item $M\geq 5$: primary obs./exp. $= 7.22$, random percentile $= 98.5$, outside-period percentile $= 100.0$.
\end{itemize}
Therefore, the active interval is not only elevated relative to the long-term mean; it also ranks among the most extreme comparable windows in the available 2020--2026 record.

\subsection{Are the anomalies localized or part of a broader regional pulse?}
The results favor a \textbf{localized-to-subregional rate pulse} rather than a spatially uniform regional increase. Several lines of evidence support this interpretation.

First, region-level full-period anomaly measures differ strongly among subregions. The \texttt{m1\_m3\_local\_union} shows the highest excess, the M2-related regions are also elevated, and the \texttt{control\_region} is much weaker. For the full active period:
\begin{itemize}[leftmargin=1.5em]
    \item \texttt{m1\_m3\_local\_union}: obs./exp. $= 10.51$, 11.14, 12.35, and 12.52 for $M\geq 3$, $M\geq 4$, $M\geq 5$, and $M\geq 6$, with percentile ranks from 98.87 to 100.0.
    \item \texttt{m2\_near\_field}: obs./exp. $= 4.82$, 5.14, 6.35, and 6.12, with percentile ranks from 95.71 to 100.0.
    \item \texttt{m2\_outer\_band}: obs./exp. $= 3.38$, 2.99, 2.79, and 6.64, with percentile ranks from 98.63 to 100.0.
    \item \texttt{control\_region}: obs./exp. $= 0.97$, 1.36, 1.90, and 1.22, with percentile ranks from 59.52 to 88.17 except for a sparse $M\geq 6$ cell.
\end{itemize}
Second, the control-resampled results show strong spatial decay. Many $>60$ km bins are near zero or unstable, while shallow 0--30 km bins in the M1--M3 local union, the common region, and selected M2-related zones remain strongly elevated. This combination indicates that the anomaly is not domain-wide and homogeneous.

Third, the spatial anomaly mapping from Task 01 is described as patchy rather than uniform, with strongest positive clusters near the M1--M3 and M2-related areas and mixed or negative cells elsewhere. Together, these outputs indicate a broader active interval with important regional coherence, but they do not support an interpretation of evenly elevated seismicity across the entire study domain.

\subsection{Which regions and depth ranges show the strongest anomalies?}
The strongest anomaly is the shallow 0--30 km portion of the \texttt{m1\_m3\_local\_union}. This is the most stable conclusion across the primary evidence matrix and the control-resampled analysis.

Representative Task 01 matrix values for \texttt{m1\_m3\_local\_union | 0--30 km} are:
\begin{itemize}[leftmargin=1.5em]
    \item $M\geq 3$: observed 637, expected 82.77, obs./exp. $= 7.70$;
    \item $M\geq 4$: observed 139, expected 15.79, obs./exp. $= 8.80$;
    \item $M\geq 5$: observed 44, expected 4.56, obs./exp. $= 9.64$.
\end{itemize}
The control summary then identifies the same cells as the top anomalies after bootstrap/random-window resampling:
\begin{itemize}[leftmargin=1.5em]
    \item \texttt{m1\_m3\_local\_union | 0--30 km | $M\geq 4$}: random ratio mean $= 94.10$, random percentile $= 99.1$, outside percentile $= 100.0$;
    \item \texttt{m1\_m3\_local\_union | 0--30 km | $M\geq 5$}: random ratio mean $= 58.97$, random percentile $= 99.1$, outside percentile $= 100.0$;
    \item \texttt{m1\_m3\_local\_union | 0--30 km | $M\geq 3$}: random ratio mean $= 32.49$, random percentile $= 98.9$, outside percentile $= 100.0$.
\end{itemize}
These cells are shown visually in Figures~\ref{fig:ratioheat} and \ref{fig:percheat}.

\begin{figure}[H]
    \centering
    \includegraphics[width=0.82\textwidth]{<CASE_ROOT>/run/02_4_bg_rate_correction/exp_run/outputs/02_aomori_bootstrap_controls/figures/01_bootstrap_control_ratio_heatmap.png}
    \caption{Control-resampled anomaly ratio heatmap for selected thresholds. The dominant hotspot is the shallow \texttt{m1\_m3\_local\_union}, with additional elevated bins in the shallow common region and selected M2-related zones.}
    \label{fig:ratioheat}
\end{figure}

\begin{figure}[H]
    \centering
    \includegraphics[width=0.82\textwidth]{<CASE_ROOT>/run/02_4_bg_rate_correction/exp_run/outputs/02_aomori_bootstrap_controls/figures/02_bootstrap_control_percentile_heatmap.png}
    \caption{Control-resampled percentile heatmap. The main target bins cluster near the 99th--100th percentile, while many far-field and deep bins are weak, zero-valued, or unstable because of sparse counts.}
    \label{fig:percheat}
\end{figure}

Beyond the M1--M3 local union, the shallow common region and selected M2-related bins are also strong. In particular:
\begin{itemize}[leftmargin=1.5em]
    \item \texttt{common\_region | 0--30 km | $M\geq 5$}: primary obs./exp. $= 7.22$, random percentile $= 98.5$;
    \item \texttt{m2\_near\_field | 30--60 km | $M\geq 4$}: primary obs./exp. $= 5.37$, random percentile $= 98.2$;
    \item \texttt{m2\_near\_field | 30--60 km | $M\geq 5$}: primary obs./exp. $= 7.31$, random percentile $= 100.0$;
    \item \texttt{m2\_outer\_band | 0--30 km | $M\geq 4$}: primary obs./exp. $= 5.05$, random percentile $= 98.9$;
    \item \texttt{m2\_outer\_band | 0--30 km | $M\geq 3$}: primary obs./exp. $= 5.53$, random percentile $= 98.9$.
\end{itemize}
Thus, the strongest signal is shallow, but some 30--60 km M2-related bins are also statistically notable.

\subsection{Does the M1--M3 local region remain anomalous after background correction?}
Yes. This is the clearest single conclusion of the analysis. The M1--M3 local union remains strongly anomalous after:
\begin{itemize}[leftmargin=1.5em]
    \item filtering the long-term raw catalog to the same common spatial domain,
    \item normalizing by long-term regional background,
    \item stratifying by depth, and
    \item applying bootstrap/random-window controls based on the long-term record.
\end{itemize}
The full-period percentile results already place this region near the top of the background distribution for every robust threshold. The control-resampled outputs then show that the shallow 0--30 km M1--M3 cells remain among the strongest anomalies in the entire matrix. This supports a firm statistical conclusion: the M1--M3 local region cannot be reduced to a minor byproduct of a weak domain-wide uplift.

\subsection{Can M2 outer-band activity be explained by broader background-rate changes?}
Only partly. The evidence indicates that the M2 outer band participates in a broader regional active interval, but it also shows localized excess beyond what is seen in the control region.

For the full active period, \texttt{m2\_outer\_band} has percentile ranks near 98.6--100.0 for $M\geq 3$ through $M\geq 6$, with obs./exp. ratios of 3.38, 2.99, 2.79, and 6.64. In the primary evidence matrix, the shallow 0--30 km outer-band values are consistently elevated across thresholds, with obs./exp. values of 4.21 for $M\geq 1.2$, 4.67 for $M\geq 2$, 5.53 for $M\geq 3$, 4.39 for $M\geq 4$, and 4.26 for $M\geq 5$.

The control results strengthen this interpretation selectively rather than uniformly. The strongest supported outer-band bins are shallow 0--30 km at $M\geq 3$ and $M\geq 4$, while deeper and far-field bins are much weaker or sparse. Meanwhile, the control region remains much less elevated. Therefore, the best-supported statistical interpretation is that the M2 outer band is \emph{not fully explained} by a generic regional background-rate rise. Instead, it appears to combine broader regional activation with its own localized enhancement, especially in the shallow depth range.

\subsection{Depth structure of the anomaly}
The anomaly is dominated by shallow seismicity. Task 01 reports that the 0--30 km depth bin carries the strongest rate peaks, with secondary contribution from 30--60 km and little support for a strong $>60$ km anomaly. Representative values include:
\begin{itemize}[leftmargin=1.5em]
    \item \texttt{common\_region | 0--30 km | $M\geq 5$}: obs./exp. $= 7.22$;
    \item \texttt{common\_region | 30--60 km | $M\geq 5$}: obs./exp. $= 2.74$;
    \item \texttt{common\_region | $>60$ km | $M\geq 5$}: obs./exp. $= 1.00$.
\end{itemize}
For M2-related regions, some 30--60 km bins are statistically strong, but the very deep bins remain weak or unstable. The control heatmaps also show many $>60$ km bins at 0 or near 0, reinforcing the interpretation that the principal anomaly is shallow to intermediate depth, not deep.

\subsection{Temporal structure: uniform rise or burst-like sequence?}
The active period is burst-like rather than uniform. Task 01 identifies three main episodes aligned with M1, M2, and M3:
\begin{itemize}[leftmargin=1.5em]
    \item an initial strong burst around M1 in mid-November 2025,
    \item the largest burst around M2 in December 2025,
    \item a renewed major burst around M3 in late April 2026.
\end{itemize}
The sequence therefore contains multiple distinct pulses within a broader elevated interval. The regional decomposition indicates partial synchrony, but not perfect regional coherence: the M1--M3 local union dominates the November and late-April phases, while the M2 outer band dominates the December pulse. This supports a segmented view of the activity rather than a single homogeneous rate increase through time.

\subsection{Robustness to control-window choice}
The control comparison using all long-term random windows and windows outside the active dates does not materially change the main conclusions. Figure~\ref{fig:scatter} shows that the two control definitions are closely aligned, with the strongest anomalies remaining strong under both schemes.

\begin{figure}[H]
    \centering
    \includegraphics[width=0.72\textwidth]{<CASE_ROOT>/run/02_4_bg_rate_correction/exp_run/outputs/02_aomori_bootstrap_controls/figures/03_random_vs_outside_control_scatter.png}
    \caption{Comparison of anomaly ratios using all long-term random windows versus windows outside the active dates. The close alignment indicates that the main anomaly conclusions are not an artifact of control-window definition.}
    \label{fig:scatter}
\end{figure}

\section{Integrated Interpretation}
The combined evidence supports the following interpretation.

\begin{enumerate}[leftmargin=1.5em]
    \item \textbf{The 2025--2026 Aomori active period is exceptional relative to the 2020--2026 background within the common analysis region.} This conclusion is supported by both direct observed-versus-expected comparisons and heavy random-window/bootstrap controls at the most reliable thresholds.
    \item \textbf{The strongest anomalies are localized, especially in the shallow M1--M3 local union.} This is the dominant result across the evidence matrix, percentile analysis, and control resampling.
    \item \textbf{The anomaly is not purely local.} The shallow common region also shows strong excess, and selected M2-related regions are statistically elevated. This indicates a broader regional active interval.
    \item \textbf{The anomaly is not spatially uniform.} The weak control-region response and weak or unstable deep/far-field bins argue against a homogeneous region-wide pulse.
    \item \textbf{The M2 outer band is statistically interesting.} Its shallow bins exceed background expectations and are stronger than the control region, so a simple explanation by broad background change alone is not adequate.
    \item \textbf{The anomaly is primarily shallow.} Intermediate-depth involvement exists in selected bins, but there is little robust evidence for a major deep anomaly.
    \item \textbf{The temporal evolution is burst-like.} The sequence contains multiple distinct pulses associated with M1, M2, and M3 rather than a single flat elevation.
\end{enumerate}

\section{Limitations and Reporting Constraints}
The conclusions above are scientifically useful but should be read with several explicit limitations.

\begin{itemize}[leftmargin=1.5em]
    \item \textbf{Active-period file substitution.} The requested active-period filename was missing. The analysis used the resolved relocated catalog at \url{<CASE_ROOT>/data/catalog/Snet_catalog_relocate_250930_260501.csv}. This does not invalidate the workflow, but it should be disclosed for reproducibility.
    \item \textbf{Magnitude completeness uncertainty.} Completeness was estimated by maximum curvature, with $M_c \approx 1.15$ for the long-term common region and $M_c \approx 1.45$ for the active common region. Therefore, low-magnitude cross-catalog comparisons are less secure than the main $M\geq 3$ to $M\geq 5$ results.
    \item \textbf{Sparse-bin inflation.} Some large obs./exp. or control-ratio values arise from very small expected counts, especially in high-magnitude or depth-partitioned bins. Such values are better interpreted as rarity indicators than as precise effect-size estimates.
    \item \textbf{Depth consistency.} The relocated catalog is shallower and underrepresents deeper seismicity relative to the raw catalog. Consequently, negative or weak findings below 60 km should remain descriptive.
    \item \textbf{Finite background window.} The reference period spans 2020-01-01 to 2026-05-22. This is adequate for the requested task but does not capture possible decade-scale nonstationarity.
    \item \textbf{No mechanism inference.} These are catalog-level statistical results only. They do not demonstrate triggering, slow slip, fluid migration, stress transfer, or fault interaction.
\end{itemize}

\section{Key Evidence Files}
The principal machine-readable outputs used in this report are listed below for traceability.

\begin{longtable}{>{\raggedright\arraybackslash}p{0.28\textwidth}>{\raggedright\arraybackslash}p{0.67\textwidth}}
\toprule
Artifact & Absolute path \\
\midrule
\endfirsthead
\toprule
Artifact & Absolute path \\
\midrule
\endhead
Common-region footprint figure & \url{<CASE_ROOT>/run/02_4_bg_rate_correction/exp_run/outputs/01_aomori_background_rate_analysis/figures/01_catalog_footprints_common_region.png} \\
Matched-event difference figure & \url{<CASE_ROOT>/run/02_4_bg_rate_correction/exp_run/outputs/01_aomori_background_rate_analysis/figures/03_matched_event_differences.png} \\
Catalog QA summary & \url{<CASE_ROOT>/run/02_4_bg_rate_correction/exp_run/outputs/01_aomori_background_rate_analysis/tables/catalog_qa_summary.csv} \\
Common-region retention & \url{<CASE_ROOT>/run/02_4_bg_rate_correction/exp_run/outputs/01_aomori_background_rate_analysis/tables/common_region_retention.csv} \\
Magnitude completeness summary & \url{<CASE_ROOT>/run/02_4_bg_rate_correction/exp_run/outputs/01_aomori_background_rate_analysis/tables/magnitude_completeness_summary.csv} \\
Long-term background rates & \url{<CASE_ROOT>/run/02_4_bg_rate_correction/exp_run/outputs/01_aomori_background_rate_analysis/tables/long_term_background_rates.csv} \\
Active-period percentile summary & \url{<CASE_ROOT>/run/02_4_bg_rate_correction/exp_run/outputs/01_aomori_background_rate_analysis/tables/active_period_percentile_vs_background.csv} \\
Primary evidence matrix & \url{<CASE_ROOT>/run/02_4_bg_rate_correction/exp_run/outputs/01_aomori_background_rate_analysis/tables/background_rate_evidence_matrix.csv} \\
Control evidence matrix & \url{<CASE_ROOT>/run/02_4_bg_rate_correction/exp_run/outputs/02_aomori_bootstrap_controls/tables/background_rate_evidence_matrix_with_controls.csv} \\
Bootstrap/random-window controls & \url{<CASE_ROOT>/run/02_4_bg_rate_correction/exp_run/outputs/02_aomori_bootstrap_controls/tables/bootstrap_random_window_controls.csv} \\
Control ratio heatmap & \url{<CASE_ROOT>/run/02_4_bg_rate_correction/exp_run/outputs/02_aomori_bootstrap_controls/figures/01_bootstrap_control_ratio_heatmap.png} \\
Control percentile heatmap & \url{<CASE_ROOT>/run/02_4_bg_rate_correction/exp_run/outputs/02_aomori_bootstrap_controls/figures/02_bootstrap_control_percentile_heatmap.png} \\
Control scatter plot & \url{<CASE_ROOT>/run/02_4_bg_rate_correction/exp_run/outputs/02_aomori_bootstrap_controls/figures/03_random_vs_outside_control_scatter.png} \\
Bootstrap/control text summary & \url{<CASE_ROOT>/run/02_4_bg_rate_correction/exp_run/outputs/02_aomori_bootstrap_controls/summary_report.txt} \\
\bottomrule
\end{longtable}

\section{Conclusions}
The statistical answer to the main question is clear: \textbf{the 2025--2026 Aomori active period is exceptional relative to the 2020--2026 background within the common analysis region}. The strongest evidence comes from the harmonized common-region comparison and is reinforced by bootstrap/random-window controls at $M\geq 3$, $M\geq 4$, and $M\geq 5$.

The most prominent anomaly is the shallow 0--30 km portion of the \texttt{m1\_m3\_local\_union}, which remains extreme after regional correction and control resampling. Selected M2-related bins are also strongly elevated, especially the shallow \texttt{m2\_outer\_band} and intermediate-depth \texttt{m2\_near\_field}. The control region is much weaker, and many deep bins are weak or unstable. Accordingly, the anomaly is best described as a \textbf{localized-to-subregional rate pulse} rather than a spatially uniform regional increase.

For later physical follow-up, the most statistically interesting observations are:
\begin{enumerate}[leftmargin=1.5em]
    \item strong common-region exceptionality across robust thresholds $M\geq 3$--$M\geq 5$ and even at $M\geq 6$ in the primary comparison;
    \item persistent, shallow, and very strong anomaly in the \texttt{m1\_m3\_local\_union};
    \item shallow excess in the \texttt{m2\_outer\_band} beyond the control-region response;
    \item selective but notable 30--60 km enhancement in parts of the M2 near field; and
    \item a burst-like temporal structure with distinct M1, M2, and M3 pulses rather than a single uniform increase.
\end{enumerate}
These findings justify targeted physical investigation, but they should remain separated from mechanism claims. The present report establishes statistical anomaly support, not causation.

\end{document}
